\title{Byte Latent Transformer: Patches Scale Better Than Tokens}

\begin{abstract}
We present the Byte Latent Transformer ({\textsc BLT}), a novel byte-level LLM framework that achieves, for the first time, performance on par with token-based LLMs at scale, alongside notable gains in inference speed and resilience. 
{\textsc BLT} processes input by encoding sequences of bytes into adaptively sized patches, which become the foundational computational units. Patch boundaries are determined according to the predicted entropy of subsequent bytes, enabling greater computational effort and capacity to be focused on segments of higher informational complexity. We conduct the first flop-constrained scaling analysis of byte-level LLMs, extending up to models with 8B parameters and trained on 4T bytes. Our experiments establish that it is practical to train large-scale models directly on raw byte sequences, eliminating the need for pre-defined vocabularies. Efficiency during both training and inference is boosted by the model’s adaptive use of longer patches in predictable contexts, and we observe enhanced qualitative results in tasks requiring reasoning and rare case generalization. In summary, when constrained to equivalent inference costs, {\textsc BLT} demonstrates clearly superior scaling characteristics compared to token-based architectures, as it is able to expand both patch length and model size concurrently.
\end{abstract}

\section{Introduction}
\label{section:intro}

We present the Byte Latent Transformer~(\textbf{BLT}), an architecture that operates without explicit tokenization by directly leveraging raw byte inputs. For the first time, this approach demonstrates parity with conventional tokenization-based systems at scale, while offering notable gains in both computational efficiency and resilience to various perturbations (\S\ref{sec:robustness}). 
In contrast, most large language models (llms) today are trained in an almost fully end-to-end manner but still rely on a preliminary, heuristic tokenization process, which converts sequences of bytes into a predefined set of tokens. This tokenization introduces compression biases that result in several limitations, including vulnerability to input domain or modality changes~\citep{dagangetting}, increased sensitivity to noisy inputs~(\S\ref{sec:robustness}), insufficient orthographic representation~\citep{edman2024cute}, and disparities across different languages~\citep{liang2023xlm,petrov2024language,limisiewicz2024myte}.

Historically, tokenization has played a critical role because training large language models on raw bytes incurs substantial computational overhead given the extended sequence lengths involved~\citep{xue2022byt5}. Earlier studies have attempted to address this issue by leveraging more computationally efficient forms of self-attention~\citep{el2020characterbert, clark2022canine} or by utilizing architectures that do not rely on attention mechanisms~\citep{wang2024mambabyte}~(\S\ref{section:related_work}). Nevertheless, such approaches mainly benefit the training of \textit{smaller models}. When scaling up, the computational burden of a Transformer chiefly stems from the substantial feed-forward network layers that process every byte, rather than from the attention components.

To optimize the distribution of computational resources, we introduce a dynamic and adaptable approach for organizing bytes into \textit{patches}~(\S\ref{section:patching}), alongside a novel model architecture that integrates both byte-level and patch-level data. In contrast to tokenization, {\textsc BLT} operates without a pre-defined patch vocabulary. Instead, it employs efficient, trainable encoder and decoder components to translate arbitrary byte groupings into corresponding latent patch embeddings. Our experiments demonstrate that this framework achieves a \textit{greater} efficiency in compute utilization compared to models reliant on tokenization.

Tokenization-based language models distribute computational resources uniformly across all tokens. This approach prioritizes convenience over optimal performance, as the process of inducing tokens relies on compression strategies that do not always align with the actual difficulty of making predictions. The core premise of our design is that computational effort should be adaptively directed to the areas that demand it most. For instance, predicting the ending letters of a word is generally a straightforward, low-entropy operation compared to determining the initial word of a new sentence, which requires more substantial computation and is often handled by large transformers. This principle is embodied in {\textsc BLT}'s architecture~(\S\ref{section:architecture}), which features three transformer blocks: two compact byte-level \textit{local models} and a sizable global \textit{latent transformer}~(\autoref{fig:arch_high_level}).
In order to organize bytes into patches—and thus control dynamic compute allocation—{\textsc BLT} uses the entropy of next-byte predictions to segment data, resulting in groups of bytes that share a similar density of information according to their contextual predictive complexity.

We introduce the first study that investigates the scaling of byte-level models while constraining by flop count, extending to models containing up to 8B parameters and exposed to 4T bytes of training data. This demonstrates the feasibility of end-to-end training at scale directly on bytes, obviating the need for fixed-vocabulary tokenization. Our findings reveal that {\textsc BLT} achieves flop-controlled training performance\footnote{Model computational requirements are assessed by enumerating all Floating Point OPerations (flops).} on par with Llama 3, yet requires up to 50\% fewer flops during inference~(\S\ref{section:scaling}).

Furthermore, we provide evidence that utilizing raw bytes leads to notable gains in handling infrequent or outlier data distributions. Compared to models that rely on tokenization, {\textsc BLT} demonstrates increased resilience to input noise and improved understanding at the character level, as exhibited in tasks probing orthographic properties, phonological patterns, and translation in low-resource settings~(\S\ref{sec:robustness}).

Additionally, {\textsc BLT} enables the concurrent scaling of both model and patch sizes while keeping inference-time flops fixed. Employing larger patch sizes typically yields computational savings, which can in turn be redirected to enlarge the global latent transformer component, since its application frequency is reduced. Our inference flop-controlled scaling experiments~(\autoref{fig:fixed_inference_scaling}) reveal marked improvements in scaling dynamics relative to approaches based on tokenization.

To conclude, our work offers the following key contributions: 1) We present {\textsc BLT}, a byte-level latent language model design that adaptively distributes computational resources to enhance flop efficiency; 2) We demonstrate flop-controlled training performance on par with Llama 3 up to the 8B parameter scale, while enabling up to 50\% improvements in flop efficiency at the expense of only modest decreases in evaluation scores; 3) The {\textsc BLT} framework introduces a novel approach for scaling language models, allowing model size increases while holding inference computational cost constant; and 4) We provide evidence that {\textsc BLT} models are more resilient to noisy inputs and capture sub-word properties of the data that are overlooked by token-based language models. All code for training and inference with {\textsc BLT} is publicly available at~\url{https://github.com/facebookresearch/blt}.

\section{Patching: From Individual Bytes to Groups of Bytes}
\label{section:patching}

Dividing a byte sequence into \textit{patches} enables {\textsc BLT} to flexibly distribute computational resources in response to contextual information. Figure~\ref{fig:patching_types} illustrates multiple approaches for partitioning bytes into patches. Mathematically, a patching function $f_p$ breaks down a byte sequence $\pmb{x}=\{x_i,|i=1,\ldots n\}$ of length $n$ into $m < n$ patches $\pmb{p}=\{p_j|j=1,\ldots,m\}$, where each $x_i$ is assigned a value from \{0,1\}—with 1 signifying the initiation of a new patch. For both token-level and patch-level architectures, the main Transformer’s processing cost is largely dictated by the total computation steps it performs. In the case of {\textsc BLT}, this count is defined by the number of patches derived from the chosen patching function. Accordingly, the typical number of bytes in each patch—referred to as \textit{patch size}—plays a crucial role in quantifying the computational cost during both model training and inference for any selected patching function~(\S\ref{section:flops}).

Subsequently, we detail three distinct patching functions: fixed-size patching where each patch contains a uniform number of bytes~(\S\ref{section:static-patch}), whitespace-based segmentation~(\S\ref{section:white-patch}), and adaptive patching leveraging entropy signals from a compact byte language model~(\S\ref{section:dyn-patch}). Additionally, we explore incremental patch construction and highlight the distinctions between tokenization and patching~(\S\ref{section:bpe-patch}).

\subsection{Strided Patching Every K Bytes}
\label{section:static-patch}

A commonly used method for organizing bytes involves segmenting them into fixed-size patches of length $k$, as described in MegaByte~\citep{yu2023megabyte}. This fixed-step approach is simple to apply during both training and inference, facilitates easy adjustment of the average patch size, and allows for straightforward management of computational requirements. Nevertheless, this strategy has notable drawbacks. Primarily, the allocation of compute resources is not adaptive to the informational content of the data: for example, a transformer layer $j$ might be underutilized when generating low-information content like whitespace, or conversely, may fall short for regions densely packed with information, such as mathematical expressions. Additionally, this method often yields patch boundaries that are inconsistent and lack contextual relevance, leading to scenarios where identical byte sequences—such as the same word—may be partitioned in disparate ways.

\subsection{Space Patching}
\label{section:white-patch}

\citet{slagle2024spacebyte} introduces a straightforward yet impactful enhancement to strided patching, wherein new patches are initiated at any space-like byte\footnote{
    Space-like bytes encompass any bytes that do not represent Latin characters, digits, or utf-8 continuation bytes. Furthermore, each patch is required to have at least one non space-like byte. }—positions commonly corresponding to linguistic unit boundaries in many languages. Under this "Space patching" approach, a latent transformer operation (with increased computational cost) is dedicated to each word, guaranteeing a consistent patching strategy for words across various sequences and ensuring that additional computational resources are devoted to the more challenging predictions that typically occur just after spaces. For illustration, generating the initial byte of the answer to ``Who composed the Magic Flute?\uline{\hspace{.6em}}'' represents a greater challenge than generating subsequent bytes after the ``M,'' given that the first letter dramatically narrows the field of probable responses, thus rendering the completion ``Mozart'' easier to predict once the initial byte is known. Nonetheless, space patching exhibits limitations in handling certain languages and domains, and crucially, it lacks the flexibility to adjust patch sizes. In what follows, we present a novel patching approach inspired by the observation that word-initial bytes tend to pose the greatest prediction difficulty, while also affording an intuitive way to modulate patch size.

\subsection{Entropy Patching: Using Next-Byte Entropies from a Small Byte LM}
\label{section:dyn-patch}

Instead of depending on heuristic rules like whitespace, we adopt a data-driven strategy to detect next-byte predictions that exhibit significant uncertainty. We propose \textit{entropy patching}, a method that utilizes entropy measurements to determine the boundaries of patches.

We train a small byte-level auto-regressive language model on the training data for {\textsc BLT} and compute next byte entropies under the LM distribution $p_{e}$ over the byte vocabulary $\mathcal{V}$:

\begin{equation}
H(x_i) = \sum_{v \in \mathcal{V}} p_{e}(x_i=v|\pmb{x}_{<i}) \log p_{e}(x_i=v|\pmb{x}_{<i})
\end{equation}

We explore two techniques for determining patch boundaries based on entropy values $H(x_i)$. The initial method selects positions where the entropy exceeds a fixed global threshold, as demonstrated in~\autoref{fig:patching}. The alternative strategy flags points where the entropy is significantly higher compared to the immediately preceding value. This latter approach can be viewed as detecting instances where the general trend of monotonically decreasing entropy within a patch is interrupted.

\begin{align*}
    \text{Global Constraint}\hspace{1cm}H(x_t)& > \theta_g\\
    \text{Approx. Monotonic Constraint}\hspace{1cm}H(x_t)& - H(x_{t-1}) > \theta_r
\end{align*}

Patch boundaries are determined through an efficient preprocessing phase that takes place during data loading. This approach contrasts with~\citet{nawrot-etal-2023-efficient}, in which a classifier is trained to recognize entropy-driven patch boundaries. Within our experiments~(\S\ref{section:experiments}), we evaluate and compare these two techniques for differentiating low-entropy from high-entropy bytes.

\subsection{The Byte-Pair Encoding (BPE) Tokenizer and Incremental Patching}
\label{section:incremental-patching}
\label{section:bpe-patch}

Numerous contemporary llms, such as our baseline model Llama 3, utilize subword tokenizers like bpe~\citep{gage1994new, sennrich-etal-2016-neural}. Here, we define ``tokens'' as byte clusters selected from a \textit{finite} vocabulary set before training, in contrast to ``patches,'' which are sequences grouped dynamically and are not constrained to a predetermined vocabulary. One notable distinction is that tokens, unlike patches, prevent the model from having explicit access to the original byte-level features.

An important advancement offered by {\textsc BLT} compared to tokenization-based approaches lies in its reconfiguration of the balance between vocabulary size and computational demands. In conventional language models, augmenting vocabulary size leads to larger average token lengths, thereby reducing the number of decoding steps required by the model; however, this simultaneously results in a corresponding increase in the output dimensionality of the model’s projection layer. Consequently, tokenization-based methods face significant constraints in altering token size and the associated inference overhead. As an illustration, Llama 3 elevates the average token size from 3.7 to 4.4 bytes, but this improvement comes at the expense of a fourfold expansion in the embedding table size relative to Llama 2.

During generation, {\textsc BLT} must determine at each step of the byte sequence whether it is at a patch boundary, as this choice dictates whether the Latent Transformer is activated for additional computation. This evaluation must be made without access to yet-ungenerated portions of the sequence. Consequently, the approach for dividing the byte sequence into patches cannot depend on knowledge of subsequent bytes. Formally, a patching scheme $f_p$ exhibits the property of incremental patching if:
$$f_p(\pmb{x}_{<i}) = f_p(\pmb{x})_{<i}$$
bpe does not possess this incremental property, since a given prefix may be tokenized in multiple ways based on how the sequence continues, thus violating the condition above\footnote{Incorporating an explicit delimiter token to signal patch boundaries can render bpe incremental, but at the cost of elongating the byte sequence.}.

\section{{\textsc BLT} Architecture}
\label{section:architecture}
The {\textsc BLT} framework comprises a sizable global autoregressive language model functioning over patch-level representations. In addition, it integrates two more compact local modules: one is responsible for converting sequences of bytes into patch representations, while the other reconstructs bytes from these patch representations~(\autoref{fig:arch_high_level}).

\subsection{Latent Global Transformer Model}

\textit{The Latent Global Transformer} refers to an autoregressive transformer, denoted as $\mathcal{G}$, composed of $l_{\mathcal{G}}$ layers. This model processes a sequence of latent input patch representations, $p_j$, and produces corresponding output patch representations, $o_j$. In our notation, $j$ is used as the index for patches, while $i$ indexes bytes. The transformer operates with a block-causal attention mask~\citep{dubey2024llama}, which ensures that attention is limited to the present and preceding patches within each document. Notably, the global model accounts for the majority of floating-point operations both in pre-training and inference phases. Consequently, strategically selecting when to apply this model enables dynamic regulation of computational resources based on the complexity of various input and output regions.

\subsection{Local Encoder}
\label{section:local}

The \textit{Local Encoder Model}, represented as $\mathcal{E}$, is constructed as a compact transformer with $l_{\mathcal{E}} << l_{\mathcal{G}}$ layers, designed to rapidly convert a series of input bytes $b_i$ into highly informative patch embeddings, $p_j$. Unlike standard transformers, it incorporates a cross-attention mechanism subsequent to each transformer layer; this modification enables pooling from byte-level to patch-level representations (see~\autoref{fig:crossattn}).

Initially, bytes $b_i$ from the input are mapped via a lookup embedding matrix of dimension $\mathbb{R}^{256 \times h_{\mathcal{E}}}$ to generate vectors $x_i$. These initial representations may be optionally enriched by the inclusion of hash-embeddings, as explained in~\S\ref{sec:hash-emb}.

Sequences of alternating transformer and cross-attention layers~(\S\ref{sec:enc-cross-attn}) then refine the representations into patch embeddings $p_i$, which are subsequently utilized by the global transformer, $\mathcal{G}$. Within each transformer layer, a \textit{local block causal} attention masking strategy is applied: every byte can access up to $w_{\mathcal{E}}$ earlier bytes, which may cross patch boundaries but are constrained within document edges. The upcoming subsections elaborate on the embedding approach and specifics of the cross-attention module.

\subsubsection{Encoder Hash n-gram Embeddings}
\label{sec:ngram-emb}
\label{sec:hash-emb}
An essential factor in developing strong and meaningful representations for each position $i$ involves leveraging context from prior bytes. {\textsc BLT} approaches this challenge by encoding both the current byte $b_i$ on its own and within the context of a byte-level n-gram. Specifically, for every position $i$, we begin by generating byte-grams

\begin{align}
    g_{i,n}=\{b_{i-n + 1},\ldots, b_i\}
\end{align}

for all byte indices $i$ and for each $n$ ranging from three to eight.\footnote{
We disregard byte-grams of length $n$ or greater when $i < n$. }

Next, we employ hash $n$-gram embeddings that project every byte $n$-gram through a hash function onto an index within an embedding matrix $E_{n}^{hash}$, where the matrix has a predetermined dimensionality, for each $n \in \{3,4,5,6,7,8\}$~\citep{bai2010learning}.
The embedding resulting from this process is subsequently summed with the original byte embedding, after which it undergoes normalization prior to being supplied as input to the local encoder architecture. The following augmented embedding is computed:

\begin{align}
e_i &= x_i + \sum_{n = 3,...,8}E_{n}^{hash}(\text{Hash}(g_{i,n})) \\
\text{where, Hash}(g_{i,n}) &= \text{RollPolyHash}(g_{i,n}) \% |E_{n}^{hash}|
\end{align}

We scale $e_i$ by dividing it by the total number of $n$-gram sizes incremented by one, utilizing RollPolyHash as specified in Appendix~\ref{appendix:rollpolyhash}. Section~\ref{section:ablations} presents an ablation study examining how varying $n$-gram hash embedding parameters—including different choices of $n$ and embedding table dimensions—affect scaling law patterns when the number of floating-point operations is controlled. Beyond hash-based $n$-gram embeddings, we also investigated $n$-gram embeddings constructed based on frequency; further information regarding these experiments can be found in Appendix~\ref{appendix:ngrams}.

\subsubsection{Encoder Multi-Headed Cross-Attention}
\label{sec:enc-cross-attn}
Our approach largely mirrors the input cross-attention mechanism used in the Perceiver model~\citep{jaegle2021perceiver}, with a key distinction: instead of relying on a predetermined set of latent representations, we use latent vectors that reflect variable-length patch representations~(\autoref{fig:crossattn}), each of which only interacts with the bytes from its associated patch. For each patch $p_j$, the corresponding query vector is derived by pooling together the byte embeddings specific to $p_j$, which is subsequently passed through a linear transformation, $\mathcal{E}_{C} \in \mathbb{R}^{h_{\mathcal{E}} \times (h_{\mathcal{E}} \times U_{\mathcal{E}})}$, where $U_{\mathcal{E}}$ denotes the number of encoder cross-attention heads. Specifically, if $f_{\text{bytes}}(p_j)$ represents the bytes comprising patch $p_j$, then the computation proceeds as follows:

\begin{align}
P_{0,j} &= \mathcal{E}_{C}(f_\text{bytes}((p_j)), f~ \text{is a pooling function} \\
P_l &= P_{l-1} + W_o\left(\text{softmax}\left(\frac{QK^T}{\sqrt{d_k}}\right)V\right) \\
\text{where } Q_j &= W_q(P_{l-1,j}), K_i = W_k(h_{l-1,i}), V_i = W_v(h_{l-1,i}) \\
h_l &= \text{Encoder-Transformer-Layer}_l(h_{l-1})
\end{align}

Here, $P \in \mathbb{R}^{n_p \times h_{\mathcal{G}}}$ denotes the representations of $n_p$ patches to be input to the global model. These representations are initially formed by aggregating the byte-level embeddings $e_i$ that correspond to each patch $p_j$. The matrices $W_q$, $W_k$, $W_v$, and $W_o$ are used for projecting queries, keys, values, and the output respectively, with the keys and values generated by projecting byte-level representations $h_i$ from the prior layer (or $e_i$ in the first layer). Our approach uses a custom masking scheme for patches: each query $Q_j$ can only attend to the keys and values associated with its respective patch $j$. Due to the use of multi-head attention over $Q, K$, and $V$, and given that patch-level representations often have greater dimensionality ($h_{\mathcal{G}}$) than $h_{\mathcal{E}}$, we represent $P_l$ as a set of heads each with dimension $h_{\mathcal{E}}$ when performing cross-attention, then concatenate these results to construct a vector of size $h_{\mathcal{G}}$. Furthermore, queries, keys, and values are each normalized using a pre-LayerNorm step, and no positional encoding is incorporated into this cross-attention mechanism. To complete the module, a residual connection is added that wraps the cross-attention block.

\subsection{Local Decoder} 

The local decoder $\mathcal{D}$, akin to its encoder counterpart, consists of a compact transformer architecture utilizing $l_{\mathcal{D}}$ layers, where $l_{\mathcal{D}}$ is substantially fewer than $l_{\mathcal{G}}$. Its role is to reconstruct the sequence of raw bytes $y_i$ from the corresponding global patch representations $o_j$. In performing this decoding, the local decoder predicts each raw byte based on the context of previously generated bytes, using as input the encoded hidden states derived from the local encoder for the byte sequence. The architecture alternates between $l_{\mathcal{D}}$ layers of cross-attention and standard transformer blocks. Within the decoder, cross-attention is employed prior to the transformer block to transform patch representations into byte-level representations, after which the local decoder’s transformer layer processes the generated byte sequence.

\subsubsection{Decoder Multi-headed Cross-Attention} 

In the decoder cross-attention, the roles of the queries and key/values are interchanged i.e. the byte-representations are now the queries, and the patch representations are now the key/values. The initial byte-representations for the cross-attention are initialized as the byte embeddings from the last encoder layer i.e. $h_{l_{\mathcal{E}}}$. The subsequent byte-representations for layer $l$, $d_{l,i}$ are computed as:

\begin{align}
D_0 &= h_{l_{\mathcal{E}}} \\
B_l &= D_{l-1} + W_o\left(\text{softmax}\left(\frac{QK^T}{\sqrt{d_k}}\right)V\right), \\
  &\text{where } Q_i = W_q(d_{l-1,i}), K_i = W_k(\mathcal{D}_C(o_j)), V_i = W_v(\mathcal{D}_C(o_j)) \\
  D_l &= \text{Decoder-Transformer-layer}_l(B_l)
\end{align}

Here, as before, $W_k$ and $W_v$ denote the key and value projection matrices applied via the linear transformation and split function $\mathcal{D}_C$, targeting the ultimate patch representations $o_j$ produced by the global model. The matrix $W_q$ performs a query projection on byte-level representations $d_{l-1}$ from the preceding decoder transformer layer (or $h_{l_{\mathcal{E}}}$ in the case of the initial layer). The output projection matrix is $W_o$. As a result, the matrix $B$ has dimensions $\mathbb{R}^{h_{\mathcal{D}} \times n_b}$, where $n_b$ stands for the total number of bytes produced in the output. Subsequent decoder states, $D_l$, are derived through a decoder transformer layer acting upon the cross-attention block output $B$. Similar to the design in the local encoder cross-attention, our implementation leverages multi-head attention, employs pre LayerNorms, omits positional encodings, and incorporates a residual connection enveloping the cross-attention module.

\section{Experimental Setup}
\label{section:experiments}
We have constructed a set of rigorously controlled experiments to directly assess {\textsc BLT} against models built on tokenization, being especially careful to avoid conferring an unfair advantage to {\textsc BLT} that might arise from exposure to longer sequence lengths.

\subsection{Pre-training Datasets} The models evaluated in this study are pre-trained on two distinct datasets: (1) The Llama 2 dataset~\citep{touvron2023llama}, consisting of 2 trillion tokens sourced from a broad range of publicly accessible materials, which undergo substantial cleaning and filtering processes to enhance data quality; and (2) {\textsc BLT}-1T: a newly constructed dataset comprising 1 trillion tokens from several public origins, which additionally integrates a portion of the pre-training data made available by Datacomp-LM~\citep{li2024datacomp}. The Llama 2 dataset is primarily utilized for scaling law investigations, focusing on the optimal token count as identified by~\cite{dubey2024llama}, to inform architectural decisions relevant to {\textsc BLT}. In contrast, the {\textsc BLT}-1T dataset is employed for end-to-end pre-training, enabling direct evaluation alongside Llama 3 on various downstream benchmarks. It is noteworthy that both datasets are strictly devoid of any information derived from Meta products or services. Additionally, baseline experiments involving tokenizer-based models utilize the Llama 3 tokenizer, which supports a 128K-token vocabulary, as this configuration provided consistently superior baseline outcomes in our assessments compared to the Llama 2 tokenizer.

\subsection{Entropy Model}
For our experiments, the entropy model is implemented as a byte-level language model that is trained using the same data distribution as the complete {\textsc BLT} model. Unless specified otherwise, our standard configuration employs a transformer consisting of 100 million parameters, structured over 14 layers with a hidden size of 512, and equipped with a sliding window attention span covering 512 bytes. All other hyperparameters mirror those utilized for both local and global transformer variants. Additional investigations, as outlined in \autoref{section:ablations}, examine variations in model capacity, input window size, and architectural choices. Notably, if the model’s receptive field is sufficiently constrained, it becomes feasible to represent the trained entropy model as a compact lookup table.

\subsection{Entropy Threshold and Equalizing Context Length} For models employing entropy-driven patching methods, we determine a patching threshold that yields a specified mean \textit{patch size} across the pretraining data mixture. In the context of {\textsc BLT}, as opposed to tokenization approaches, the \textit{patch size} is not constrained and may be set arbitrarily, which has considerable impact on the effective context window available to the model. To ensure a fair comparison and prevent models with larger patch sizes from benefiting disproportionately from increased context, we fix the total expected number of bytes per batch. Accordingly, when the patch size is increased, the sequence length is proportionally reduced. For Llama 2 data, we implement an 8k byte context window, whereas for the {\textsc BLT}-1T corpus, this is expanded to an average of 16k bytes, all while retaining a constant batch size of 16M bytes on average.

Although the mean batch size does not change, dynamic patching approaches can cause variation in the proportion of bytes per patch during batch loading. To optimize efficiency, our {\textsc BLT} training implementation groups patches into batches in such a way that padding is minimized within the resource-intensive latent transformer module. This guarantees a consistent number of patches across all batches. To prevent memory usage from spiking due to exceptionally large patches, we pad—or, when necessary, truncate—byte sequences to fixed lengths of 12k and 24k bytes for Llama 2 and {\textsc BLT}-1T datasets, respectively, throughout the training process.

\subsection{Entropy Model Context}
\label{section:entropy-context}
Our empirical observations indicate that applying entropy patching results in increasingly larger patch sizes within structured content, such as multiple choice formats that are frequently repetitive (refer to the MMLU case in \autoref{fig:mmlu_patching}). These changes arise due to decreased entropy in the repeated segments present in the entropy model's context. Therefore, for the large-scale application of {\textsc BLT}-Entropy with a patch size of 4.5, we reinitialize the entropy context by inserting new lines and utilize the approximate monotonicity constraint, which is less vulnerable to the phenomenon of "entropy drift" triggered by varying context lengths. Importantly, this adjustment impacts only the calculation method for entropies, while the approach for determining the entropy threshold value remains consistent.

\subsection{FLOPs Estimation} 
\label{section:flops}

Our methodology for calculating transformer flops aligns closely with that outlined in Chinchilla~\citep{hoffmann2022training}, accounting for flop usage in the feed-forward modules, the qkvo projections within the self-attention mechanism, and operations for both attention computation and output projection. One important distinction in our approach is the assumption that the input embedding layer leverages an efficient lookup method rather than relying on dense matrix multiplication, rendering this step effectively flop-free. Consistent with prior literature, we approximate the computational demand of the backward pass to be double that of the forward pass.

To compute flops \textit{per byte} for {\textsc BLT} models, we add up the flops for the local encoder transformer, the global latent transformer, and the local decoder transformer, together with the cross attention blocks in the encoder and the decoder:

\begin{align}
    \text{FL}_{\text{{\textsc BLT}}} &= \text{Transf. FL}(h_{\mathcal{G}}, l_{\mathcal{G}}, m=n_{ctx}/n_p, V=0) / n_p\\
    &+ \text{Transf. FL}(h_{\mathcal{E}}, l_{\mathcal{E}}, m=w_{\mathcal{E}}, V=0) \\
    &+ \text{Transf. FL}(h_{\mathcal{D}}, l_{\mathcal{D}}, m=w_{\mathcal{D}}, V=256) \\ 
    &+ \text{Cross Attn. FL}(h_{\mathcal{E}}, l_{\mathcal{E}}, m=n_p, r=n_p/k) \cdot k/n_p \\ 
    &+ \text{Cross Attn. FL}(h_{\mathcal{D}}, l_{\mathcal{D}}, m=k, r=k/n_p)
\end{align}

In this context, $n_{ctx}$ represents the sequence length measured in bytes, while $n_p$ denotes the patch size. The variable $r$ specifies the proportion between the number of queries and the number of keys/values, and $k$ reflects the ratio of the patch dimension to the byte dimension; this indicates how many local model segments are concatenated to construct the full global model representation (for example, $k=2$ is shown in \autoref{fig:crossattn}). The parameter $V$ indicates the output projection's vocabulary size, which is utilized exclusively within the local decoder. The value of $m$, denoting the context length used for attention mechanisms, differs based on whether the module processes the sequence at the byte or patch level. As a result, we adapt the calculation of attention-related floating point operations (flops) for each module accordingly. Detailed formulas for the computation of Transformer-FLOPs and Cross-Attention FLOPs are included in Appendix~\ref{section:flops-appendix}.

\subsection{Bits-Per-Byte Estimation} Perplexity only makes sense in the context of a fixed tokenizer as it is a measure of the uncertainty for each token. When comparing byte and token-level models, following previous work~\citep{xue2022byt5, yu2023megabyte, wang2024mambabyte}, we instead report Bits-Per-Byte (BPB), a tokenizer independent version of perplexity. Specifically:

\begin{align}
    \text{BPB}(x)&= \frac{\mathcal{L}_{CE}(\pmb{x})}{\ln(2)\cdot n_{\text{bytes}}}
\end{align}

In this formulation, the aggregated cross-entropy loss indicating the uncertainty for $\pmb{x}$ is divided by both the quantity of bytes in $\pmb{x}$ and a scaling constant for normalization.

\subsection{Transformer Architecture Hyperparameters}  
All transformer layers in {\textsc BLT}, covering both the local and global modules, adopt hyperparameter settings broadly aligned with the Llama~3 architecture~\citep{dubey2024llama}. Specifically, the feed-forward components utilize the SwiGLU activation~\citep{swiglu}. For positional encoding, rotary positional embeddings (RoPE)~\citep{RoPe} are applied solely within the self-attention mechanisms, with $\theta=500000$ as recommended by~\citet{xiong2024effective}. Layer normalization is achieved with RMSNorm~\citep{RMSNorm}. In all self-attention layers employing standard fixed attention patterns—such as \textit{block causal} or \textit{fixed-window block causal}—Flash attention~\citep{dao2022flashattention} is used. These fixed-width masks are configured to a window size of 512. Due to the dynamic, patch-specific masking inherent to our cross-attention layers, we adopt Flex Attention\footnote{\url{https://pytorch.org/blog/flexattention}}; this facilitates efficient fused implementations and yields substantial reductions in training time.

\subsection{{\textsc BLT}-Specific Hyperparameters} To evaluate the performance of {\textsc BLT} models, we perform experiments in two main areas: scaling behavior and downstream task performance, considering model sizes of 400M, 1B, 2B, 4B, and 8B parameters. Detailed architectural hyperparameters for each model variant can be found in Appendix~Table~\ref{tab:arch}. Queries for the initial cross-attention layer within the local encoder are initialized using max-pooling. Across all {\textsc BLT} models, we employ $500,000$ hashes with a single hash function, utilizing n-grams of lengths 3 to 8. All models are trained with a learning rate set to $4e-4$. The equivalence in optimal learning rates between token-based and {\textsc BLT} models is established through a hyperparameter sweep within the range $1e-3$ to $1e-4$ at the 400M and 1B scales, where the shared value proved most effective. For experiments evaluating scaling with Llama-2 data, we adhere to batch sizes as recommended by~\cite{dubey2024llama} or their bytewise equivalents. Model optimization is carried out using AdamW~\citep{adamw} with parameters $\beta_1 = 0.9$, $\beta_2 = 0.95$, and $\epsilon=10^{-8}$. We employ a linear learning rate warm-up over 2000 steps followed by a cosine decay to zero, weight decay set to 0.1, and apply global gradient clipping at a maximum norm of 1.0.

\section{Scaling Trends}
\label{section:scaling}

We provide a comprehensive overview of scaling behaviors in byte-level models to guide the continued development of {\textsc BLT} models. Our investigation into scaling addresses the gaps present in earlier byte-level modeling studies in several key respects: (a) We analyze patterns within the context of compute-efficient training regimes, (b) We develop and assess corresponding 8B models utilizing substantial training data volumes (up to 1T tokens or 4T bytes) with downstream evaluation, and (c) We examine how scaling dynamics operate under fixed inference-cost scenarios. A subsequent section delves into particular benefits associated with processing byte-sequences.

\subsection{Parameter Matched Compute Optimal Scaling Trends}
\label{section:parameter-matched-scaling}
To investigate compute-optimal scaling, we utilize the Llama 2 dataset to train multiple bpe and {\textsc BLT} models spanning four parameter regimes, from 1B to 8B parameters. We evaluate model performance by plotting the total number of training flops against language modeling accuracy on a representative sample drawn from the training data mixture. The bpe models are optimized according to the model-to-data ratio established as optimal by Llama~3~\citep{dubey2024llama}, ensuring that they follow the theoretical compute-optimal regime intended to maximize performance given a fixed computational budget~\citep{hoffmann2022training}. This approach establishes a strong reference point. Each bpe configuration is paired with a {\textsc BLT} counterpart, for which we employ a Latent Transformer identical in size and underlying architecture to the equivalent bpe Transformer, training both variants on the same dataset.

As shown in~\autoref{fig:scaling-compute-opt} (right), {\textsc BLT} models consistently perform at least on par with, and often surpass, their bpe-based equivalents, with this pattern persisting as both model scale and compute increase. To our knowledge, {\textsc BLT} represents the first byte-level Transformer framework to demonstrate scaling behaviors comparable to those of BPE architectures within compute-optimal settings. This finding supports our hypothesis that the parameter-to-compute ratio optimal for bpe models also extends to {\textsc BLT}, or deviates only minimally.

Achieving bpe scaling behavior necessitates both enhancements in architecture and the incorporation of dynamic patching. In ~\autoref{fig:scaling-compute-opt} (left), we present a comparison of models employing space patching and Llama 3. Our estimation of SpaceByte \citep{slagle2024spacebyte} relies on {\textsc BLT} space-patching, but excludes both n-gram embeddings and cross-attention modules. While SpaceByte achieves improvements relative to Megabyte, its performance does not approach that of Llama 3. The right panel of \autoref{fig:scaling-compute-opt} demonstrates the collective impact of architectural refinements and dynamic patching. Notably, at scale, {\textsc BLT} models rival the leading tokenizer-based models, including Llama 3.

Our results further reveal that the selection of tokenizer notably impacts performance in tokenizer-based models; specifically, models utilizing the Llama-3 tokenizer demonstrate superior results compared to those employing the Llama-2 tokenizer, despite being trained on identical datasets.

In summary, the {\textsc BLT} architecture demonstrates performance situated between Llama 2 and Llama 3 when much larger patch sizes are employed. While the bpe tokenizers for Llama 2 and 3 exhibit mean token lengths of 3.7 and 4.4 bytes, respectively, {\textsc BLT} achieves comparable scaling dynamics even when the average patch size is increased to 6 or 8 bytes. Since the number of inference flops is inversely related to the average patch size, selecting an 8-byte patch size can reduce inference flops by up to 50\%. Moreover, the use of larger patch sizes appears advantageous as both model scale and data volume increase. For example, {\textsc BLT} configured with an 8-byte patch size initially underperforms relative to bpe Llama 2 at the 1B scale, but surpasses it by the time model size reaches 7B parameters. This trend implies that increasing patch sizes may yield further improvements at even greater scales, hinting that employing still larger patch sizes could be viable as the capacity of models and the amount of training compute expand.

\subsection{Beyond Compute Optimal Task Evaluations}
\label{section:evals}
To further investigate scaling behavior, we train an 8B {\textsc BLT} model past the compute-optimal regime, utilizing the {\textsc BLT}-1T dataset—a significantly larger and higher quality corpus—and examine its capabilities across multiple widely used classification and generation benchmarks. For our evaluation, we include tasks that assess common sense reasoning, general world knowledge, and programming ability:
\paragraph{Classification tasks} encompass ARC-Easy (0-shot)~\citep{Eval_ARC}, Arc-Challenge (0-shot)~\citep{Eval_ARC}, HellaSwag (0-shot)~\citep{Eval_hellaswag}, PIQA (0-shot)~\citep{Eval_piqa}, and MMLU (5-shot)~\citep{hendrycks2020measuring}. Performance is measured using a prompt scoring approach, in which the model's likelihood assignments to possible answer choices are aggregated, and mean accuracy is reported.
\paragraph{Coding related generation tasks:}  To assess Python code synthesis capabilities, we provide pass@1 results on MBPP (3-shot)~\citep{Eval_MBPP} and HumanEval (0-shot)~\citep{Eval_HumanEval}.

In \autoref{tab:evals}, a comparison is presented among three distinct models, each trained on the {\textsc BLT}-1T dataset: a Llama 3 tokenizer-based bpe model,\footnote{We opt for the Llama 3 tokenizer, which features a 128k vocabulary, as it yields better performance relative to Llama 2's 32k vocabulary.} as well as two different variants of the {\textsc BLT} model. The first {\textsc BLT} variant adopts a space-patching approach ({\textsc BLT}-Space), whereas the second leverages an entropy-driven patching method ({\textsc BLT}-Entropy), applying an approximate monotonicity constraint and resetting the entropy model’s context on encountering new lines (as elaborated in~\autoref{section:entropy-context}). Training for all three models is conducted under comparable flop budgets. For {\textsc BLT}-Entropy, we further apply an inference-time modification, lowering the entropy threshold from 0.6 to 0.1. This adjustment is observed to enhance performance across tasks, although it does so at the expense of an increased number of inference steps.

The {\textsc BLT}-Entropy model achieves superior results compared to the Llama 3 model on 4 of the 7 evaluated tasks, despite both models being trained on an equivalent amount of data measured in bytes. This performance gain is likely attributable to two primary factors: (1) more efficient utilization of computational resources during training through dynamic patching, and (2) the explicit incorporation of byte-level data rather than relying solely on token representation.

Conversely, {\textsc BLT}-Space lags behind the Llama 3 tokenizer in all but one evaluation task. Nevertheless, its average patch size of 6 bytes leads to a notable decrease in inference flops. By contrast, models based on bpe and entropy-patching techniques display similar average patch sizes, each around 4.5 bytes on the combined training set. Given an equivalent training budget, the model with the larger patch size is able to process 30\% more data compared to the other two; this increased coverage may, however, drive {\textsc BLT} further from achieving compute-optimality.

\subsection{Patches Scale Better Than Tokens} The {\textsc BLT} framework permits simultaneous scaling of both model and patch sizes without exceeding the original training and inference computational cost or increasing the volume of training data. The ability to expand patch size without bound is a distinctive property of patch-centric models, enabling them to circumvent the efficiency limitations inherent in fixed-vocabulary token-based approaches, as outlined in Section~\ref{section:bpe-patch}. Utilizing larger patches reduces computational demands, thereby freeing up resources that can be redirected to increase the capacity of the global latent transformer, since its invocation frequency is correspondingly diminished.

In this work, we perform a fixed inference scaling analysis to examine whether scaling up model size while reducing the number of processing steps per input, particularly with larger patch sizes, yields better performance relative to deploying smaller models with more steps. Using models initialized with 400 million and 3.6 billion parameters along with the Llama 2 tokenizer, we identify flop-equivalent variants employing the Llama 3 tokenizer, as well as {\textsc BLT}-Entropy models configured with average patch sizes of 6 and 8 bytes on our training data mixture (refer to~\autoref{tab:fixed-inf-params} for detailed model configurations). For models with a patch size of 8, we adopt 3 encoder layers, in contrast to 1 layer for other setups. Across these configurations, we train each model variant using a range of predetermined training flop budgets.

As illustrated in \autoref{fig:fixed_inference_scaling}, {\textsc BLT} models demonstrate superior scaling behavior compared to tokenization-based counterparts across both inference flop categories. BPE architectures initially deliver stronger performance when the training budget is limited, but they are quickly outperformed by {\textsc BLT} models as the training compute approaches and exceeds the compute-optimal point. In practical settings, allocating greater resources to the single pretraining phase can yield models with substantially improved performance for a given inference constraint. This phenomenon is exemplified by the class of 8B models, such as Llama 3.1, which has undergone pretraining on a volume of data two orders of magnitude larger than the compute-optimal amount for its size.

The point at which {\textsc BLT} surpasses token-based models has moved marginally closer to the compute-optimal threshold as we transition to larger flop class models, shifting from 3x to approximately 2.5x the optimal compute budget. Moreover, models employing the larger patch size of 8 exhibit a steeper improvement trajectory within the higher flop regime, surpassing other variants at an earlier stage. As explored in Section~\ref{section:parameter-matched-scaling}, it is observed that models with increased patch sizes tend to approach the performance of BPE models as the scale grows. This phenomenon can be partially explained by the relatively reduced proportion of total flops allocated to the byte-level Encoder and Decoder modules, which display slower scaling characteristics compared to the Latent Transformer. For instance, when expanding overall model parameters by 20-fold (from 400M to 8B), the parameters dedicated to the local structure of {\textsc BLT} only experience a twofold increase. This is a significant detail, since enlarging patch size principally impacts flops within the patch Latent Transformer component while leaving the byte-level modules largely unaffected. Consequently, this explains the observation that the {\textsc BLT}-Entropy with a patch size of 8 rose from 1.6x to 1.7x the parameter count of the Llama 2 model as model scale increased.

In conclusion, our investigation into scaling with respect to patch length reveals that the {\textsc BLT} patch-based model architecture benefits from coordinated scaling of both patch size and overall model capacity. This positive scaling behavior not only holds but appears to strengthen as the model size increases further.

\section{Byte Modeling Improves Robustness}
Additionally, we evaluate the robustness of {\textsc BLT} against token-based counterparts that do not incorporate byte-level data. We also introduce a methodology for converting existing pretrained token-based models to operate at the byte level.
\label{sec:robustness}

\subsection{Character-Level Tasks}

An initial impetus for developing byte-level models stemmed from their ability to handle input corrupted by noise at the byte level, as well as from their sensitivity to sub-token elements—a capability often lacking in models that rely on tokenization. To assess these attributes, we conduct supplementary experiments on benchmarks specifically designed to gauge resistance to input noise and to test character-level comprehension in both English and multiple other languages, encompassing not only alphabetic characters but also digits and phonemes. The outcomes of these assessments are summarized in Table~\ref{tab:char_tasks}.

\paragraph{Noisy Data} To assess model robustness, we generate altered versions of the benchmark classification datasets referenced in Section~\ref{section:evals}. Specifically, we compare how tokenizer-based models and {\textsc BLT} handle noisy inputs by implementing five different character-level perturbations: (a) \textit{AntSpeak}: renders the entire text as uppercase letters, separated by spaces; (b) \textit{Drop}: randomly eliminates 10\% of the text’s characters; (c) \textit{RandomCase}: randomly assigns half of the characters to uppercase and half to lowercase; (d) \textit{Repeat}: randomly selects 20\% of the characters and repeats them up to four times; (e) \textit{UpperCase}: capitalizes every character in the sequence. For evaluation, these noising operations are individually applied to the prompt, the completion, or both, and performance metrics are averaged across these conditions. Table \ref{tab:char_tasks} displays the corresponding results for noised HellaSwag~\citep{Eval_hellaswag}. The findings reveal that {\textsc BLT} demonstrates significantly greater resilience to textual noise than its tokenizer-based counterparts, exhibiting an average performance lead of 8 points over an equivalently trained tokenizer-based model, and even surpassing the Llama 3.1 model, which was trained on a substantially larger dataset.

\paragraph{Phonology - Grapheme-to-Phoneme (G2P)} We evaluate the proficiency of {\textsc BLT} in converting a sequence of graphemes, which are the written symbols constituting a word, into its corresponding phonemic transcription. Table \ref{tab:char_tasks} reports G2P task results under a 5-shot evaluation framework utilizing Phonology Bench~\citep{suvarna2024phonologybench}. Our findings demonstrate that {\textsc BLT} surpasses the Llama 3 1T model utilizing the tokenizer-based baseline for this particular task.

\paragraph{CUTE}
To evaluate character-level comprehension, we test {\textsc BLT} using the CUTE benchmark~\citep{edman2024cute}. This suite of tasks is organized into three principal areas: compositional understanding, orthographic similarity recognition, and sequence manipulation. The CUTE benchmark presents substantial difficulties for most models relying on tokenizers, which generally retain some awareness of token spelling but often fail to exploit this for effective text manipulation. As shown in Table~\ref{tab:char_tasks}, {\textsc BLT}-Entropy exceeds the performance of BPE Llama 3 models by over 25 points on these tasks. Our model is particularly adept at character manipulation, attaining 99.9\% accuracy on both spelling-oriented tasks. Such pronounced gains, achieved even though {\textsc BLT} was trained on only one-sixteenth of the data used for Llama 3.1, suggest that BPE-based models struggle to internalize character-level patterns. Figure \ref{fig:cute_examples} presents representative cases in which the Llama 3 tokenizer model encounters difficulty, while our {\textsc BLT} model excels. BPE models only surpass in the word deletion and word insertion tasks; these forms of word-level manipulation may present unique challenges to byte-level models, though the gap remains modest, and compositional learning from characters up to words may ultimately prove more manageable than the inverse. For all tasks, we preserve the same evaluation protocol and utilize the original Huggingface prompts. It is possible that more effective prompt engineering could further assist BPE models.

\paragraph{Low Resource Machine Translation} We assess the performance of {\textsc BLT} on translation tasks involving both directions across six prominent language families and twenty-one lower resource languages—spanning diverse scripts—utilizing the FLORES-101 benchmark~\citep{goyal2022flores}. SentencePiece BLEU scores are presented in Table~\ref{tab:flores}. The evaluation shows that {\textsc BLT} delivers superior results compared to a counterpart trained with the Llama 3 tokenizer, offering a consistent 2-point improvement in translations into English and a 0.5-point increase for translations from English. While performance with major language pairs is on par with, or slightly surpasses, Llama 3, {\textsc BLT} exhibits a notable edge across many pairings within lower-resource language families. These findings highlight the utility of byte-level modeling in enabling robust generalization to rare and complex byte sequences.

\subsection{Training {\textsc BLT} via Llama 3 Initialization}
\label{sec:distillation}
We investigate a procedure in which {\textsc BLT} models are able to capitalize on established pre-trained tokenizer-based architectures to facilitate more rapid and stable training progress. In particular, this is implemented by seeding {\textsc BLT}'s global transformer parameters with those from a pre-trained Llama 3.1 model. During subsequent training, we adopt a learning rate for the global transformer that is one-tenth that used for both the local encoder and local decoder, adhering to the recommended step count for Llama 3. Performance outcomes, alongside a baseline {\textsc BLT}, are shown in Table~\ref{tab:distill}. The results clearly demonstrate that initializing {\textsc BLT} with weights from Llama 3.1 yields pronounced improvements over both the standalone Llama 3 and a naively trained {\textsc BLT} baseline, with all models trained using equivalent FLOPs. Furthermore, this method even surpasses our {\textsc BLT}-Entropy variant (Table~\ref{tab:evals}), which utilized a much larger 1T-token dataset, on the MMLU benchmark. This indicates that leveraging Llama 3.1 in this fashion can be a highly efficient means of reducing computational resources required for training.

This approach can be interpreted as converting models reliant on tokenizers into models that operate without them, thereby adapting a pre-trained LLaMA 3.1 model into a {\textsc BLT} variant. For a thorough comparison, we list the baseline LLaMA 3.1 model—trained with 15T tokens—in Table \ref{tab:distill}, juxtaposed with the performance of the LLaMA-derived {\textsc BLT}. While our adapted model demonstrates a moderate decrease in performance on MMLU and HumanEval, its results drop more considerably on several other benchmarks. These outcomes highlight the necessity for further research to better harness pre-trained models, particularly through refinement of data mixture strategies and careful tuning of additional hyperparameters.

\section{Ablations and Discussion}
\label{section:ablations}
Here, we examine a series of ablation studies that support our decisions regarding the architecture of {\textsc BLT}, as well as the selection of patching methods and hyper-parameter configurations for the 8B parameter {\textsc BLT} model trained on the {\textsc BLT}-1T dataset.

\paragraph{Entropy Model Hyper-parameters} In order to examine how changes in the size of the entropy model and the context window affect scaling, we train a range of byte-level entropy transformer models varying in size from 1 million to 100 million parameters, using context window lengths spanning 64 to 512. Utilizing our $400m$ and $1b$ {\textsc BLT} models trained on the Llama-2 dataset, we construct and display training flop scaling law curves of bpb in \autoref{fig:entropy_ablation}. Our results indicate that increased model size and larger context windows generally enhance scaling behavior, though improvements taper off beyond models of 50 million parameters.

\paragraph{Types of Patching}

We evaluate the impact of each of the four patching methods described in Section~\ref{section:patching}, namely: 1) Strided Patching employing strides of 4 and 6, 2) Patching at whitespace locations, 3) BPE Tokenizer patching utilizing the Llama 3 tokenizer, and 4) Entropy-based patching guided by a compact byte-level language model.

Although dynamic patching shortens the actual sequence length, we regulate the sequence length so that each patching scheme is provided with an equivalent context length. During both training and inference, every model is exposed to an identical expected number of bytes per sequence, thereby eliminating confounding variables related to modeling longer contexts. Figure~\ref{fig:scaling-compute-opt} presents the outcomes of these ablation studies. Among the approaches, all alternative patching methods demonstrate superior performance over static patching, and notably, space patching shows performance nearly comparable to that of dynamic entropy-based patching.

\autoref{tab:evals_ablation} displays the benchmark results for {\textsc BLT} models, offering a comparative analysis of tokenizer-based approaches, space patching, and entropy-based patching when trained on the Llama 2 dataset for the optimal number of steps as outlined by~\citet{dubey2024llama}. While space patching provides a more straightforward alternative that circumvents the need to execute an entropy model during training, our findings indicate that the improvements attributed to entropy-based patching in scaling experiments~(see Section~\ref{section:scaling}) also extend to various downstream benchmark tasks.\footnote{Results for space patching reflect earlier experiments without cross-attention, though the overall trend persists even with cross-attention enabled.}

\paragraph{Cross-Attention} 
As shown in~\autoref{tab:crossattn_ablations_1b}, we examine the effect of inserting cross-attention at different stages within the encoder and decoder components of {\textsc BLT}. For the encoder, we explore three strategies for initializing the cross-attention queries: (1) assigning all global states a shared learned embedding, (2) employing a hash-based embedding derived from the patch's bytes, and (3) applying pooling operations to the encoder's hidden states for the patch bytes at the current layer.

Our results indicate that incorporating cross-attention within the \textit{decoder} yields the highest efficacy. Within the encoder, cross-attention offers modest gains, which manifest primarily when the queries are initialized through pooling. Furthermore, cross-attention proves especially beneficial for the Common-Crawl dataset and is increasingly advantageous as patch size grows.

\paragraph{n-gram Hash Embeddings} 

We experiment with n-gram hash embedding vocabulary sizes of 0, 100K, 200K, and 400K, and summarize the outcomes in Table~\ref{tab:ngram_ablations_1b}. The data indicate that the inclusion of hash embeddings leads to improvements across all domains, with especially notable gains on Wikipedia and Github (a 0.04 bpb increase compared to 0.01 bpb following 15k steps at the 8B model scale). At the 8B scale, reducing the number of hashes from 500K to 300K results in only a minor decrease in performance (around 0.001 bpb after 15k steps). These findings suggest that hash embeddings play a critical role in narrowing the performance gap between {\textsc BLT} and traditional tokenizer-based architectures. Nevertheless, beyond 300K hashes, additional increases yield minimal benefits. Furthermore, these improvements tend to supplement the impact of cross-attention, as the two strategies enhance outcomes on distinct datasets.

\paragraph{Local Model Hyperparameters} Table \ref{tab:local_layers} presents an ablation study exploring different configurations for the number of layers in both the local encoder and decoder. The results indicate that {\textsc BLT}, when utilizing hash n-gram embeddings, performs effectively with a very minimal encoder—specifically, a single layer—while benefitting from a more substantial, deeper decoder architecture.

\section{Related Work}
\label{section:related_work}

\textbf{Character-Level RNNs:} Character Language Modeling has remained a prominent benchmark since the inception of neural language models~\citep{sutskever2011generating,mikolov2012subword,graves2013generating} due to its intrinsic ability to handle previously unseen words naturally, obviating the need for explicit back-off strategies. \cite{kim2016character} introduced a framework that represents textual input solely through characters, leveraging convolutional and highway networks to encode the character sequences, which are subsequently fed into LSTM-based RNNs. This approach achieved competitive results with the then state-of-the-art RNN language models for English and delivered superior performance on morphologically complex languages, highlighting another key benefit of character-level LLMs. Focusing on morphologically complex languages such as Turkish and Russian, \cite{kenter2018byte} deployed byte-level LSTM architectures for machine comprehension, yielding results that surpassed those attained by word-based models. Similarly, \cite{zhang2015character} demonstrated that character-based convolutional networks outperformed word-level models in certain classification scenarios. Hierarchical modeling was explored by \citet{chung2019hierarchical}, who utilized LSTM networks augmented with boundary detectors at each level, thereby unveiling latent structural hierarchies in text and enhancing character-level language modeling capabilities. Distinctively, ByteNet~\cite{kalchbrenner2016neural} relied on convolutional neural network layers applied at the character level, instead of attention mechanisms, to address machine translation tasks.

\textbf{Character-Level Transformers:} The introduction of attention-based transformer architectures~\citep{vaswani2017attention} in conjunction with subword tokenization techniques~\citep{sennrich-etal-2016-neural} led to notable advancements in the performance of neural networks on both language modeling and standard benchmarks. Nonetheless, the use of word and sub-word representations inherently encodes a bias regarding the model’s abstraction granularity. In an effort to merge the advances of transformer models with encouraging outcomes from character-level language modeling, \citet{al2019character} constructed highly deep transformer networks augmented with auxiliary objectives, successfully training models that surpassed previous LSTM-based character-level language models. Despite these improvements, a considerable performance gap in comparison to word-level LLMs remained. In a similar vein, GPT-2~\citep{radfordgpt2} found that on extensive corpora such as the 1 Billion Word benchmark, language models operating at the byte level still lagged behind those employing word-level representations.

\cite{Choe2019BridgingTG} showed that transformer-based byte-level LLMs can exceed the performance of subword-level models with a similar number of parameters. However, these byte-level architectures demand significantly greater computational resources and longer training times. In a related effort, \cite{el2020characterbert} introduced CharFormer, a BERT variant which generates word representations via convolutions applied to character embeddings. This approach led to gains within the medical domain, yet incurred high computational costs. CANINE, presented in \cite{clark2022canine}, is an encoder-only model with 150 million parameters that operates over raw character sequences. Its architecture consists of a deep transformer reminiscent of our global module, paired with local transformer components and strided convolutions for input character downsampling. CANINE surpassed the token-based mBERT model on multilingual benchmarks. ByT5~\citep{xue2022byt5} extended byte-level modeling to encoder-decoder frameworks without employing patching or downsampling. While this contributed to increased robustness to input noise and allowed them to reach comparable performance to traditional token-based models using just a quarter of the data, the approach required costly attention computations at the byte level, resulting in heavy computational burdens. Modeling language at the byte level, as opposed to using subword units, leads to longer input sequences, which poses obstacles for scalable and efficient model training. Most recently, leveraging the Mamba Architecture~\citep{gu2023mamba}—capable of maintaining a fixed-size memory over long contexts—\cite{wang2024mambabyte} trained a byte-level Mamba model, also eschewing patching operations. In a controlled FLOP regime and at the 350 million parameter size, their approach outperformed byte-level transformer architectures in bits-per-byte metrics across multiple datasets.

\textbf{Patching-based approaches:} Leveraging patching techniques has proven effective in significantly reducing the computational costs—measured in floating point operations—associated with byte-level LLMs, all while potentially maintaining comparable performance. Several studies have showcased promising results with such methods on relatively modest model sizes and dataset scales. \cite{nawrot-etal-2022-hierarchical} explore the use of fixed patching via both downsampling and upsampling operations, proposing the hourglass transformer, which demonstrates superior results over other byte-level baselines at the 150M parameter scale. In a subsequent advancement, \cite{nawrot-etal-2023-efficient} introduce dynamic patching methods that include a learnable boundary-predictor trained in an end-to-end manner, a boundary-predictor with supervision from various tokenizers, and an entropy-guided patching strategy inspired by {\textsc BLT}. They report that these innovations surpass contemporary vanilla transformer baselines on language modeling benchmarks, especially for models with approximately 40M parameters and trained on around 400M tokens. \cite{lester2024training} take a different route by compressing input sequences via arithmetic coding to realize compression ratios that exceed what BPE provides; with their equal-info windows strategy, they manage to outdo byte-level models in language modeling metrics, though they still fall short compared to subword-level approaches.

Our research takes substantial influence from MegaByte~\citep{yu2023megabyte}, a decoder-only causal large language model (LLM) that relies on a fixed, static patching scheme with representation concatenation to transform bytes into patches. On the decoder side, a local module reconverts these patches into bytes. The MegaByte model has been shown to perform comparably to models leveraging tokenizers when evaluated at the 1B parameter scale over datasets comprising approximately 400B bytes. In our study, we perform thorough ablation of MegaByte in all empirical investigations, noting that the static patching approach fails to match the performance of contemporary, compute-efficient, state-of-the-art tokenizer-based models under flop-constrained conditions. We illustrate how {\textsc BLT} helps to close this performance gap. Similarly, \cite{slagle2024spacebyte} observe the same limitation with MegaByte and propose advancing the static patching strategy by including whitespace and other space-like bytes, along with the integration of a local encoder component. Their findings indicate gains over tokenized transformer architectures in certain domains, like Github and arXiv, at the 1B parameter level in compute-normalized comparisons. We incorporate experiments on this model in our work as well, and highlight that further architectural enhancements are necessary to enable byte-level models to scale and match the results of leading token-based systems like Llama 3.

\section{Limitations and Future Work}

For this study, our model training duration aligns with the optimal number of steps established for Llama~3~\citep{dubey2024llama}, serving as a guide for our architectural decisions. Nonetheless, since these scaling laws were derived from transformers utilizing BPE-level tokenization, their application to {\textsc BLT} may produce suboptimal ratios between dataset size and parameter count. Determining appropriate scaling laws specifically for {\textsc BLT}, which could yield more advantageous scaling characteristics for our design, remains an avenue for future investigation. It should also be noted that a significant portion of these experiments were carried out with models containing up to 1B parameters; as we expand to models with 8B parameters or more, there is potential for revised architectural configurations that may better exploit the advantages at increased scale.

Current transformer libraries and codebases prioritize optimization for architectures utilizing tokenizers. Although our study includes theoretically matched flop experiments and incorporates some optimized solutions (like FlexAttention) for non-standard transformer layers, the actual running time of our models may still lag behind those based on tokenizers, indicating potential for additional improvements in efficiency.

Although {\textsc BLT} employs an independently trained entropy model for patching, developing a patching model via end-to-end learning represents a promising avenue for future research. In Section \ref{sec:distillation}, we report preliminary results that suggest the viability of ``byte-ifying'' models like Llama 3—originally trained on over 10T tokens and reliant on tokenizers—by initializing the global transformer with their pretrained weights and subsequently freezing it. Continuing to pursue this line of research may lead to approaches that preserve the advantages of bytefying and even surpass the capabilities of these tokenizer-based models, all without necessitating complete retraining.

\section{Conclusion}
\label{section:conclusion}

This work introduces the Byte Latent Transformer (\textbf{BLT}), an innovative model architecture that departs from the standard reliance on predetermined token vocabularies in large language models. By implementing a dynamic and trainable mechanism for assembling bytes into patches, {\textsc BLT} enables adaptive allocation of computational effort according to the underlying data complexity. This strategy yields marked gains in computational efficiency as well as model robustness. Through an extensive series of scaling experiments, we show that {\textsc BLT} matches the performance of established tokenization-centric models such as Llama 3 for model sizes up to 8B and data scales reaching 4T bytes. Moreover, it can achieve reductions in inference floating point operations by up to 50\%, with only marginal impact on evaluation metrics. {\textsc BLT} also introduces an additional scaling axis by supporting simultaneous growth in both model size and patch size without exceeding a predetermined inference cost, a property particularly suited to typical computational environments. The direct processing of raw byte-level data enables {\textsc BLT} to better capture rare and complex data cases, thereby enhancing resilience to input noise and offering superior modeling of sub-word phenomena. Collectively, these findings highlight {\textsc BLT} as a robust and flexible substitute for conventional tokenization-based systems, furnishing a more adaptable and resource-efficient platform for language model development.

\end{document}